\scanpage{24}
% archive foliation: 23
\margin{\uline{Leray}}
\diagram{\begin{tikzcd}[row sep=small] F & \\ X\arrow[u,hook]\arrow[d] & \\ Y & \end{tikzcd}}
\unsure{application} fibrée (au sens de Serre) générale
\[
H^*(X)\Longleftarrow H^p(Y,H^q(F))
\]
\figure{\begin{tikzpicture}\draw (0,0) .. controls (0.4,-.08) and (2.2,.03) .. (2.8,0);\end{tikzpicture}}

\margin{\uline{Leray--\unsure{Čech}}}
\diagram{\begin{tikzcd}[row sep=small] X\arrow[d,two heads]\\Y\arrow[d]\\B_F\end{tikzcd}}
$F$ un groupe, fibré principal
\[
\leftarrow\left\{\begin{gathered}\unsure{\text{d'où}}\ \text{fibration}\\\text{générale}\end{gathered}\right.
\qquad
H^*(Y)\Longleftarrow H^p(B_F,\{H^q(X)\})
\]
N.B. $\{H^q(X)\}$ est un système local \deleted{\uline{\unsure{trivial}}}\inserted{défini} par les opérations de $\pi_0(F)$ sur $H^q(X)$.
\figure{\begin{tikzpicture}\draw (0,0) .. controls (0.7,.07) and (2.4,.1) .. (2.9,0);\end{tikzpicture}}

\diagram{\begin{tikzcd}[row sep=small]Y\arrow[d,two heads]\\B_F\arrow[d]\\B_{\deleted{\unsure{G}}\inserted{X}}\end{tikzcd}}
$\deleted{\unsure{G}}\inserted{X}$ un groupe

$F$ un ss-groupe
\[
\left\{\unsure{\text{d'où}}\ \text{fibration générale}\right.
\qquad
H^*(B_F)\Longleftarrow H^{\unsure{p}}(B_{\deleted{\unsure{G}}\inserted{X}},\{H^q(Y)\})
\]
N.B. $\{H^q(Y)\}$ est un système local défini par les op. de $\pi_0(X)$ sur $H^q(X/F)$ \unsure{via celles de} $\pi_0(X)$ sur $\pi_0(X/F)$
\figure{\begin{tikzpicture}\draw (0,0) -- (.14,.12) -- (.03,-.02) -- (.26,.13) -- (.18,0) -- (.4,.12) -- (.32,-.01) .. controls (.9,-.08) and (1.6,.08) .. (2.2,0) -- (3,0);\end{tikzpicture}}
\margin{\uline{Shapiro}}
Si $Y$ est discret, relation d'isomorphisme bien connue entre la coh. de $B_F$, et celle de $B_G$ à valeurs dans un coeff. \uline{induit} [à considérer surtout dans le cas où $F,G$ sont discrets]
\figure{\begin{tikzpicture}\draw (0,0) .. controls (.4,-.08) and (2.4,-.02) .. (3,.06);\end{tikzpicture}}

\margin{\uline{Hochschild--Serre} $\unsure{\parallel}$ \uline{Leray--\unsure{Čech}}}
\diagram{\begin{tikzcd}[row sep=small]B_F\arrow[d,two heads]\\B_X\arrow[d]\\B_Y\end{tikzcd}}
$X$ un groupe, $F$ un ss-groupe \uline{distingué} donc $Y$ un groupe
\[
\left\{\unsure{\text{d'où}}\ \text{fibration générale}\right.
\qquad
H^*(B_G)\Longleftarrow H^p(B_Y,H^q(B_F))
\]
